\input{../../../common/labpreamb_new.tex}

\renewcommand{\labnumber}{5}
\renewcommand{\tasks}{(5, 12, 2)}

\begin{document}

\maketitle
\thispagestyle{empty}
\newpage
\pagenumbering{arabic}
\tableofcontents

\newpage
\section{Задание}
\begin{enumerate}
    \item Создать  файл для хранения действительных чисел, вводимых с
        клавиатуры. Прочитать этот файл и вычислить максимальное
        среди отрицательных чисел.
    \item Создать связанный список для хранения целых чисел. Число
        записей неизвестно, но можно ввести число, являющееся
        признаком окончания ввода данных. Для четных номеров
        вариантов организовать очередь, а для нечетных – стек. Для
        исходного списка решить следующую задачу: вставить заданное
        число A1 после каждого элемента с нечетным номером.
    \item Для полученного списка решить следующую задачу: упорядочить
        по убыванию методом <<пузырька>> элементы списка,
        расположенные после первого положительного элемента.
\end{enumerate}

\newpage
\section{Внешняя спецификация}
\specframe{
    Лабораторная работа №\labnumber
}
\specframe{
    Задание №1
}
\specframeloop{
    Введите $k$ от $1$ до $ARRAY_SIZE$: \\
    $<k>$
}
До $1 \leq k \leq ARRAY\_SIZE$
\specframe{
    Введите $<<k>>$ строк длины не более $STRING\_SIZE$: \\
    $<S[0]>, <S[1]>, ..., <S[k-1]>$
}
При $n = 0$
\specframeif{
    Не нашлось подстрок, ограниченных точками
}
Иначе
\specframeif{
    Нашлось $<<n>>$ подстрок, ограниченных точками: \\
    $\ll substrs[0] \gg, \ll substrs[1] \gg, ..., \ll substrs[n - 1] \gg$
}
При $found.size = 0$
\specframeif{
    Не нашлось строки со скобками и цифрами
}
Иначе
\specframeif{
    Нашлась строка со скобками и цифрами: "$<<found>>$" \\
    После очистки букв не из русского алфавита: "$<<found>>$"
}

\newpage
\section{Листинг программы}

% \lstinputlisting[language=C]{main.c}

\newpage
\section{Распечатка тестов к программе и результатов}

\begin{center}
    \begin{tabular}{|c|l|l|}
        \hline
        № & Искодные данные & Результат \\\hline
        1 &
        $k = 1, S[0:0] = [\text{abba}]$ &
        \makecell[l]{
            Не нашлось подстрок, ограниченных точками \\
            Не нашлось строки со скобками и цифрами
        }\\\hline
        2 &
        $k = 2, S[0:2] = [\text{y..b..c ..d}, \text{.dog
        при(1-2)ветfe.}]$ &
        \makecell[l]{
            Нашлось 3 подстрок, ограниченных точками: \\
            b \\
            c  \\
            dog при(1-2)ветfe \\
            Нашлась строка со скобками и цифрами: \\
            "dog при(1-2)ветfe" \\
            После очистки букв не из русского алфавита: \\
            "привет"
        }\\\hline
        3 &
        $k = 1, S[0:0] = [\text{.adas.dog приветfeline.}]$ &
        \makecell[l]{
            adas \\
            dog приветfeline \\
            Не нашлось строки со скобками и цифрами
        }\\\hline
    \end{tabular}
\end{center}

\end{document}
